\documentclass[10pt,a4paper]{article}

%double si on veut une version prof et une version élève. simple si on veut une seule version.
\newcommand{\buildmode}{simple}

\InputIfFileExists{version.tex}{}{}
% Version (définie par le script)
\providecommand{\version}{eleve}

\usepackage{feuilleinterro}

%%%à modifier à chaque fois

\renewcommand{\niveau}{1G}
\renewcommand{\chapitre}{Suites, généralités}
\renewcommand{\numchap}{1}
\renewcommand{\theme}{Suites}

%%%titre "CORRECTION" au lieu de "INTERROGATION", sans les champs Nom/Prénom
\renewcommand{\interroversion}[1]{
    \noindent\textbf{\niveau}

    \vspace{-0.5cm}
    \begin{center}
        \textbf{\large CORRECTION -- \chapitre}
    \end{center}

    \vspace{0.1cm}
}

\begin{document}

% =========================================================
% PAGE 1 : VERSION A
% =========================================================

\interroversion{A}

\textbf{Exercice 1 -- Calcul avec des fractions \hfill /2}

Calculer, en détaillant les étapes, et donner le résultat sous la forme d'une fraction irréductible :
\[
A = \frac{2}{3}\times\frac{3}{4}-\frac{1}{6}.
\]

On effectue d'abord le produit, puis la soustraction :
\[
\begin{aligned}
A &= \frac{2\times3}{3\times4}-\frac{1}{6}\\
  &= \frac{6}{12}-\frac{1}{6}\\
  &= \frac{1}{2}-\frac{1}{6}\\
  &= \frac{3}{6}-\frac{1}{6}\\
  &= \frac{2}{6}=\frac{1}{3}.
\end{aligned}
\]

\textbf{Exercice 2 -- Suite définie par récurrence \hfill /4}

On considère la suite \((u_n)\) définie pour tout entier naturel \(n\) par
\[
\left\{\begin{array}{rcl}
u_0 &=& 1 \\
u_{n+1} &=& 2u_n+1
\end{array}\right.
\]

\begin{enumerate}
\item Quel est le mode de définition de cette suite ? Justifier.

Cette suite est définie par récurrence : le terme \(u_{n+1}\) est donné à partir du terme précédent \(u_n\) (et non directement en fonction de \(n\)), ainsi que la valeur du premier terme \(u_0\).

\item Calculer \(u_1\), \(u_2\) et \(u_4\).

\[
\begin{aligned}
u_1 &= 2u_0+1 = 2\times1+1 = 3\\
u_2 &= 2u_1+1 = 2\times3+1 = 7\\
u_3 &= 2u_2+1 = 2\times7+1 = 15\\
u_4 &= 2u_3+1 = 2\times15+1 = 31.
\end{aligned}
\]
\end{enumerate}

\textbf{Exercice 3 -- Sens de variation d'une suite \hfill /4}

On considère la suite \((u_n)\) définie pour tout entier naturel \(n\) par \(u_n=3n-5\).

Déterminer le sens de variation de la suite \((u_n)\) (on pourra calculer \(u_{n+1}-u_n\) en fonction de \(n\)).

\[
\begin{aligned}
u_{n+1}-u_n &= \bigl(3(n+1)-5\bigr)-\bigl(3n-5\bigr)\\
            &= 3n+3-5-3n+5\\
            &= 3.
\end{aligned}
\]
Pour tout entier naturel \(n\), \(u_{n+1}-u_n=3>0\), donc la suite \((u_n)\) est croissante.

\newpage

% =========================================================
% PAGE 2 : VERSION B
% =========================================================

\interroversion{B}

\textbf{Exercice 1 -- Calcul avec des fractions \hfill /2}

Calculer, en détaillant les étapes, et donner le résultat sous la forme d'une fraction irréductible :
\[
B = \frac{5}{6}\times\frac{2}{3}-\frac{1}{2}.
\]

On effectue d'abord le produit, puis la soustraction :
\[
\begin{aligned}
B &= \frac{5\times2}{6\times3}-\frac{1}{2}\\
  &= \frac{10}{18}-\frac{1}{2}\\
  &= \frac{5}{9}-\frac{1}{2}\\
  &= \frac{10}{18}-\frac{9}{18}\\
  &= \frac{1}{18}.
\end{aligned}
\]

\textbf{Exercice 2 -- Suite définie par récurrence \hfill /4}

On considère la suite \((u_n)\) définie pour tout entier naturel \(n\) par
\[
\left\{\begin{array}{rcl}
u_0 &=& 2 \\
u_{n+1} &=& 2u_n-1
\end{array}\right.
\]

\begin{enumerate}
\item Quel est le mode de définition de cette suite ? Justifier.

Cette suite est définie par récurrence : le terme \(u_{n+1}\) est donné à partir du terme précédent \(u_n\) (et non directement en fonction de \(n\)), ainsi que la valeur du premier terme \(u_0\).

\item Calculer \(u_1\), \(u_2\) et \(u_4\).

\[
\begin{aligned}
u_1 &= 2u_0-1 = 2\times2-1 = 3\\
u_2 &= 2u_1-1 = 2\times3-1 = 5\\
u_3 &= 2u_2-1 = 2\times5-1 = 9\\
u_4 &= 2u_3-1 = 2\times9-1 = 17.
\end{aligned}
\]
\end{enumerate}

\textbf{Exercice 3 -- Sens de variation d'une suite \hfill /4}

On considère la suite \((u_n)\) définie pour tout entier naturel \(n\) par \(u_n=2n-7\).

Déterminer le sens de variation de la suite \((u_n)\) (on pourra calculer \(u_{n+1}-u_n\) en fonction de \(n\)).

\[
\begin{aligned}
u_{n+1}-u_n &= \bigl(2(n+1)-7\bigr)-\bigl(2n-7\bigr)\\
            &= 2n+2-7-2n+7\\
            &= 2.
\end{aligned}
\]
Pour tout entier naturel \(n\), \(u_{n+1}-u_n=2>0\), donc la suite \((u_n)\) est croissante.

\newpage

% =========================================================
% PAGE 3 : VERSION C
% =========================================================

\interroversion{C}

\textbf{Exercice 1 -- Calcul avec des fractions \hfill /2}

Calculer, en détaillant les étapes, et donner le résultat sous la forme d'une fraction irréductible :
\[
C = \frac{3}{5}\times\frac{5}{6}-\frac{1}{3}.
\]

On effectue d'abord le produit, puis la soustraction :
\[
\begin{aligned}
C &= \frac{3\times5}{5\times6}-\frac{1}{3}\\
  &= \frac{15}{30}-\frac{1}{3}\\
  &= \frac{1}{2}-\frac{1}{3}\\
  &= \frac{3}{6}-\frac{2}{6}\\
  &= \frac{1}{6}.
\end{aligned}
\]

\textbf{Exercice 2 -- Suite définie par récurrence \hfill /4}

On considère la suite \((u_n)\) définie pour tout entier naturel \(n\) par
\[
\left\{\begin{array}{rcl}
u_0 &=& 0 \\
u_{n+1} &=& 3u_n+2
\end{array}\right.
\]

\begin{enumerate}
\item Quel est le mode de définition de cette suite ? Justifier.

Cette suite est définie par récurrence : le terme \(u_{n+1}\) est donné à partir du terme précédent \(u_n\) (et non directement en fonction de \(n\)), ainsi que la valeur du premier terme \(u_0\).

\item Calculer \(u_1\), \(u_2\) et \(u_4\).

\[
\begin{aligned}
u_1 &= 3u_0+2 = 3\times0+2 = 2\\
u_2 &= 3u_1+2 = 3\times2+2 = 8\\
u_3 &= 3u_2+2 = 3\times8+2 = 26\\
u_4 &= 3u_3+2 = 3\times26+2 = 80.
\end{aligned}
\]
\end{enumerate}

\textbf{Exercice 3 -- Sens de variation d'une suite \hfill /4}

On considère la suite \((u_n)\) définie pour tout entier naturel \(n\) par \(u_n=4n-3\).

Déterminer le sens de variation de la suite \((u_n)\) (on pourra calculer \(u_{n+1}-u_n\) en fonction de \(n\)).

\[
\begin{aligned}
u_{n+1}-u_n &= \bigl(4(n+1)-3\bigr)-\bigl(4n-3\bigr)\\
            &= 4n+4-3-4n+3\\
            &= 4.
\end{aligned}
\]
Pour tout entier naturel \(n\), \(u_{n+1}-u_n=4>0\), donc la suite \((u_n)\) est croissante.

\newpage

% =========================================================
% PAGE 4 : VERSION D
% =========================================================

\interroversion{D}

\textbf{Exercice 1 -- Calcul avec des fractions \hfill /2}

Calculer, en détaillant les étapes, et donner le résultat sous la forme d'une fraction irréductible :
\[
D = \frac{7}{10}\times\frac{5}{6}-\frac{1}{4}.
\]

On effectue d'abord le produit, puis la soustraction :
\[
\begin{aligned}
D &= \frac{7\times5}{10\times6}-\frac{1}{4}\\
  &= \frac{35}{60}-\frac{1}{4}\\
  &= \frac{7}{12}-\frac{1}{4}\\
  &= \frac{7}{12}-\frac{3}{12}\\
  &= \frac{4}{12}=\frac{1}{3}.
\end{aligned}
\]

\textbf{Exercice 2 -- Suite définie par récurrence \hfill /4}

On considère la suite \((u_n)\) définie pour tout entier naturel \(n\) par
\[
\left\{\begin{array}{rcl}
u_0 &=& 1 \\
u_{n+1} &=& 4u_n-2
\end{array}\right.
\]

\begin{enumerate}
\item Quel est le mode de définition de cette suite ? Justifier.

Cette suite est définie par récurrence : le terme \(u_{n+1}\) est donné à partir du terme précédent \(u_n\) (et non directement en fonction de \(n\)), ainsi que la valeur du premier terme \(u_0\).

\item Calculer \(u_1\), \(u_2\) et \(u_4\).

\[
\begin{aligned}
u_1 &= 4u_0-2 = 4\times1-2 = 2\\
u_2 &= 4u_1-2 = 4\times2-2 = 6\\
u_3 &= 4u_2-2 = 4\times6-2 = 22\\
u_4 &= 4u_3-2 = 4\times22-2 = 86.
\end{aligned}
\]
\end{enumerate}

\textbf{Exercice 3 -- Sens de variation d'une suite \hfill /4}

On considère la suite \((u_n)\) définie pour tout entier naturel \(n\) par \(u_n=5n-6\).

Déterminer le sens de variation de la suite \((u_n)\) (on pourra calculer \(u_{n+1}-u_n\) en fonction de \(n\)).

\[
\begin{aligned}
u_{n+1}-u_n &= \bigl(5(n+1)-6\bigr)-\bigl(5n-6\bigr)\\
            &= 5n+5-6-5n+6\\
            &= 5.
\end{aligned}
\]
Pour tout entier naturel \(n\), \(u_{n+1}-u_n=5>0\), donc la suite \((u_n)\) est croissante.

\end{document}
